\FloatBarrier
\section{نتایج ارزیابی \lr{v2}}
\label{sec:v2-results}

همه‌ی اعداد این بخش روی همان سری‌های آزمونی است که مدل پایه با آن‌ها سنجیده شد، و با نرخ ۵۰۰ هرتز. پس مقایسه‌ی دو مدل منصفانه است. آموزش هر شرکت‌کننده روی \lr{GTX 1660 Ti} بین ۱۷ تا ۳۱ دقیقه طول کشید (جدول~\ref{tab:v2-train}).

\subsection{دقت \lr{K}، \lr{E} و \lr{F}}

جدول~\ref{tab:v2-subj} خلاصه‌ی اصلی نتایج است.

\paragraph{\lr{K} از مدل پایه بهتر است.} میانگین \lr{VAF} \lr{K} \lr{۷۵٫۸} درصد و میانگین $r$ آن \lr{۰٫۸۷۲} است؛ مدل پایه \lr{۷۲٫۲} درصد و \lr{۰٫۸۵۱} داشت. این بهبود در هر ۱۲ شرکت‌کننده دیده می‌شود، نه فقط در میانگین. اندازه‌ی بهبود از \lr{۰٫۳} واحد (\lr{P5}) تا \lr{۵٫۵} واحد (\lr{P4}) است و میانه‌اش حدود ۴ واحد. \lr{K} حتی از مدل فقط‌سینماتیکِ \lr{residual experiment} (\lr{۷۵٫۵} درصد و \lr{۰٫۸۶۹}) هم کمی بهتر است. پس اولین معیار پذیرش \lr{v2} برآورده شد.

\begin{table}[htbp]
\centering
\small
\caption{نتایج \lr{v2} روی سری‌های آزمون در برابر مدل پایه. چهار ستون اول \lr{VAF} (درصد) و چهار ستون بعد \lr{r}، هر دو میانگین روی پنج عضله.}
\label{tab:v2-subj}
\begin{tabular}{lcccccccc}
\toprule
شرکت‌کننده & پایه & \lr{K} & \lr{E} & \lr{F} & پایه & \lr{K} & \lr{E} & \lr{F} \\
\midrule
\lr{P1} & \lr{۷۹٫۰} & \lr{۸۳٫۶} & \lr{۳۳٫۹} & \lr{۸۳٫۸} & \lr{۰٫۸۹۱} & \lr{۰٫۹۱۴} & \lr{۰٫۵۹۷} & \lr{۰٫۹۱۵} \\
\lr{P2} & \lr{۷۵٫۲} & \lr{۷۸٫۸} & \lr{۵۸٫۳} & \lr{۷۸٫۸} & \lr{۰٫۸۷۲} & \lr{۰٫۸۸۸} & \lr{۰٫۷۶۶} & \lr{۰٫۸۸۸} \\
\lr{P3} & \lr{۶۷٫۵} & \lr{۶۹٫۸} & \lr{۳۰٫۷} & \lr{۷۰٫۲} & \lr{۰٫۸۱۱} & \lr{۰٫۸۳۷} & \lr{۰٫۵۹۲} & \lr{۰٫۸۳۶} \\
\lr{P4} & \lr{۷۶٫۰} & \lr{۸۱٫۵} & \lr{۵۷٫۷} & \lr{۸۱٫۸} & \lr{۰٫۸۸۰} & \lr{۰٫۹۰۵} & \lr{۰٫۷۵۹} & \lr{۰٫۹۰۵} \\
\lr{P5} & \lr{۶۳٫۵} & \lr{۶۳٫۸} & \lr{۳۴٫۴} & \lr{۶۴٫۰} & \lr{۰٫۷۷۷} & \lr{۰٫۷۹۵} & \lr{۰٫۶۰۳} & \lr{۰٫۷۹۳} \\
\lr{P6} & \lr{۷۴٫۷} & \lr{۷۶٫۷} & \lr{۳۵٫۸} & \lr{۷۷٫۰} & \lr{۰٫۸۶۵} & \lr{۰٫۸۷۷} & \lr{۰٫۶۲۶} & \lr{۰٫۸۷۷} \\
\lr{P7} & \lr{۷۹٫۱} & \lr{۸۲٫۶} & \lr{۵۷٫۷} & \lr{۸۲٫۵} & \lr{۰٫۸۹۹} & \lr{۰٫۹۱۰} & \lr{۰٫۷۵۶} & \lr{۰٫۹۱۱} \\
\lr{P8} & \lr{۵۴٫۳} & \lr{۵۸٫۴} & \lr{۳۲٫۷} & \lr{۵۸٫۰} & \lr{۰٫۷۶۴} & \lr{۰٫۷۹۲} & \lr{۰٫۵۸۲} & \lr{۰٫۷۹۹} \\
\lr{P9} & \lr{۷۹٫۳} & \lr{۸۳٫۵} & \lr{۴٫۵} & \lr{۸۳٫۵} & \lr{۰٫۸۹۷} & \lr{۰٫۹۱۸} & \lr{۰٫۴۷۸} & \lr{۰٫۹۱۸} \\
\lr{P10} & \lr{۷۴٫۹} & \lr{۸۰٫۲} & \lr{۵۳٫۶} & \lr{۸۰٫۳} & \lr{۰٫۸۷۱} & \lr{۰٫۸۹۸} & \lr{۰٫۷۲۵} & \lr{۰٫۸۹۸} \\
\lr{P11} & \lr{۷۳٫۶} & \lr{۷۸٫۱} & \lr{۳۶٫۰} & \lr{۷۸٫۱} & \lr{۰٫۸۵۸} & \lr{۰٫۸۸۴} & \lr{۰٫۶۴۹} & \lr{۰٫۸۸۴} \\
\lr{P12} & \lr{۶۹٫۸} & \lr{۷۲٫۵} & \lr{۳۴٫۴} & \lr{۷۲٫۳} & \lr{۰٫۸۳۲} & \lr{۰٫۸۴۵} & \lr{۰٫۶۰۴} & \lr{۰٫۸۴۵} \\
\midrule
\textbf{میانگین} & \lr{۷۲٫۲} & \lr{۷۵٫۸} & \lr{۳۹٫۱} & \lr{۷۵٫۹} & \lr{۰٫۸۵۱} & \lr{۰٫۸۷۲} & \lr{۰٫۶۴۵} & \lr{۰٫۸۷۲} \\
\bottomrule
\end{tabular}
\end{table}


\paragraph{\lr{E} از \lr{EEG} به‌تنهایی پیش‌بینی معناداری می‌سازد.} \lr{E} که فقط \lr{EEG} می‌بیند به $r$ \lr{۰٫۶۴۵} و \lr{VAF} \lr{۳۹٫۱} درصد رسید. این از خروجی فقط‌\lr{EEG} آزمایش‌های قبلی ($r$ حدود \lr{۰٫۴۷}) بالاتر است.

آیا ممکن است این عدد تصادفی باشد؟ برای پاسخ، \lr{null test} انجام شد: بلوک‌های \lr{EEG} در زمان جابه‌جا می‌شوند تا رابطه‌اش با حرکت از بین برود، و دقت دوباره حساب می‌شود. \lr{E} واقعی در هر ۱۲ شرکت‌کننده از این حالت جابه‌جاشده بهتر بود (جدول~\ref{tab:v2-verdict}). پس \lr{EEG} به‌تنهایی اطلاعات واقعی درباره‌ی حرکت دارد.

\paragraph{\lr{F} همان \lr{K} است.} میانگین \lr{VAF} \lr{F} \lr{۷۵٫۹} و \lr{K} \lr{۷۵٫۸} است. بزرگ‌ترین قدرمطلق \lr{gate} در هر شرکت‌کننده بین \lr{۰٫۰۱} و \lr{۰٫۰۴} ماند، و اصلاحی که \lr{EEG} اضافه می‌کند حدود \lr{۰٫۲} تا ۲ درصد اندازه‌ی \lr{K} است (جدول~\ref{tab:interp-parts}). فقط در یک نفر (\lr{P4}) بهتر بودن \lr{F} نسبت به \lr{K} از \lr{null test} گذشت؛ از ۱۲ آزمون، یک نتیجه‌ی معنادار همان چیزی است که از شانس انتظار می‌رود.

این یعنی \lr{EEG} چیزی به سینماتیک اضافه نمی‌کند. ولی شاید مشکل از ماژول \lr{fusion} باشد که نتوانسته اطلاعات موجود را یاد بگیرد. بخش~\ref{sec:residual} این احتمال را جداگانه بررسی می‌کند و نشان می‌دهد این‌طور نیست.

\begin{table}[htbp]
\centering
\small
\caption{\lr{null tests}: \lr{E} در برابر \lr{EEG} جابه‌جاشده بین بلوک‌ها، و بهره‌ی \lr{F} نسبت به \lr{K}. \lr{p} هیچ‌گاه کمتر از یک‌تقسیم‌بر‌تعداد بلوک‌ها گزارش نمی‌شود. ستون آخر بزرگ‌ترین قدرمطلق \lr{gate} \lr{fusion} بین پنج عضله است.}
\label{tab:v2-verdict}
\begin{tabular}{cccccc}
\toprule
شرکت‌کننده & \lr{r} خروجی \lr{E} & \lr{p} (جابه‌جایی) & $\Delta$\lr{VAF} (\lr{F}$-$\lr{K}) & \lr{p} (جابه‌جایی) & بیشینه‌ی $|g|$ \\
\midrule
\lr{P1} & \lr{۰٫۵۹۷} & \lr{۰٫۰۱۲} & \lr{+۰٫۱۹} & \lr{۰٫۱۹} & \lr{۰٫۰۱۵} \\
\lr{P2} & \lr{۰٫۷۶۶} & \lr{۰٫۰۱۱} & \lr{−۰٫۰۲} & \lr{۰٫۸۵} & \lr{۰٫۰۱۴} \\
\lr{P3} & \lr{۰٫۵۹۲} & \lr{۰٫۰۱۳} & \lr{+۰٫۳۷} & \lr{۰٫۹۵} & \lr{۰٫۰۱۷} \\
\lr{P4} & \lr{۰٫۷۵۹} & \lr{۰٫۰۱۲} & \lr{+۰٫۳۲} & \lr{۰٫۰۱} & \lr{۰٫۰۲۶} \\
\lr{P5} & \lr{۰٫۶۰۳} & \lr{۰٫۰۱۱} & \lr{+۰٫۲۵} & \lr{۰٫۱۸} & \lr{۰٫۰۲۶} \\
\lr{P6} & \lr{۰٫۶۲۶} & \lr{۰٫۰۱۱} & \lr{+۰٫۳۵} & \lr{۰٫۲۶} & \lr{۰٫۰۲۴} \\
\lr{P7} & \lr{۰٫۷۵۶} & \lr{۰٫۰۱۱} & \lr{−۰٫۰۲} & \lr{۰٫۷۸} & \lr{۰٫۰۱۸} \\
\lr{P8} & \lr{۰٫۵۸۲} & \lr{۰٫۰۱۳} & \lr{−۰٫۳۵} & \lr{۰٫۳۴} & \lr{۰٫۰۴۲} \\
\lr{P9} & \lr{۰٫۴۷۸} & \lr{۰٫۰۱۲} & \lr{+۰٫۰۳} & \lr{۰٫۳۴} & \lr{۰٫۰۱۲} \\
\lr{P10} & \lr{۰٫۷۲۵} & \lr{۰٫۰۱۲} & \lr{+۰٫۰۶} & \lr{۰٫۴۰} & \lr{۰٫۰۱۰} \\
\lr{P11} & \lr{۰٫۶۴۹} & \lr{۰٫۰۱۱} & \lr{+۰٫۰۱} & \lr{۰٫۴۵} & \lr{۰٫۰۱۵} \\
\lr{P12} & \lr{۰٫۶۰۴} & \lr{۰٫۰۱۰} & \lr{−۰٫۲۲} & \lr{۰٫۸۲} & \lr{۰٫۰۲۱} \\
\bottomrule
\end{tabular}
\end{table}


\paragraph{کی می‌توان به \lr{F} اعتماد کرد.} \lr{F} روی \lr{K} ساخته می‌شود، پس \lr{fusion} ذخیره‌شده فقط وقتی معنا دارد که با همان \lr{K} ذخیره‌شده آموزش دیده باشد. دفترچه‌ی آموزش این شرط را بررسی می‌کند. با قاعده‌ی \lr{freezing}، بهترین دوره‌ی \lr{EEG} و \lr{fusion} همیشه بعد از \lr{freeze} شدن \lr{K} است (جدول~\ref{tab:v2-train})، پس شرط در هر ۱۲ نفر برقرار است.

جدول حکم‌هایی که این گزارش از آن استفاده می‌کند پیش از اصلاح این بررسی ذخیره شده است. به همین دلیل برای ۹ نفر برچسب «غیرقابل‌تفسیر» دارد. اعداد جدول درست‌اند؛ اگر سلول ارزیابی دوباره اجرا شود، این برچسب‌ها به «\lr{EEG} چیزی به سینماتیک اضافه نمی‌کند» تغییر می‌کنند (به‌جز \lr{P4}).

\begin{table}[htbp]
\centering
\small
\caption{روند آموزش \lr{v2}: شماره‌ی بهترین دوره و دوره‌ی \lr{freezing} هر جریان. خط تیره یعنی جریان تا پایان (۱۵۰ دوره) منجمد نشد.}
\label{tab:v2-train}
\begin{tabular}{ccccccc}
\toprule
شرکت‌کننده & دوره‌ها & بهترین \lr{K} & \lr{freezing} \lr{K} & بهترین \lr{E} & \lr{freezing} \lr{E} & زمان (دقیقه) \\
\midrule
\lr{P1} & \lr{۱۵۰} & \lr{۸۰} & \lr{۱۱۰} & \lr{۱۲۷} & -- & \lr{۲۴٫۱} \\
\lr{P2} & \lr{۱۵۰} & \lr{۱۱۶} & \lr{۱۴۶} & \lr{۱۴۷} & -- & \lr{۳۰٫۸} \\
\lr{P3} & \lr{۱۰۵} & \lr{۳۱} & \lr{۶۱} & \lr{۷۵} & \lr{۱۰۵} & \lr{۱۷٫۲} \\
\lr{P4} & \lr{۱۳۸} & \lr{۵۵} & \lr{۸۵} & \lr{۱۰۸} & \lr{۱۳۸} & \lr{۲۳٫۹} \\
\lr{P5} & \lr{۱۵۰} & \lr{۵۵} & \lr{۸۵} & \lr{۱۲۵} & -- & \lr{۲۷٫۹} \\
\lr{P6} & \lr{۱۲۲} & \lr{۴۹} & \lr{۷۹} & \lr{۹۲} & \lr{۱۲۲} & \lr{۲۳٫۱} \\
\lr{P7} & \lr{۱۰۵} & \lr{۳۸} & \lr{۶۸} & \lr{۷۵} & \lr{۱۰۵} & \lr{۲۱٫۹} \\
\lr{P8} & \lr{۱۳۸} & \lr{۲۷} & \lr{۵۷} & \lr{۱۰۸} & \lr{۱۳۸} & \lr{۲۲٫۹} \\
\lr{P9} & \lr{۱۵۰} & \lr{۷۳} & \lr{۱۰۳} & \lr{۱۲۰} & \lr{۱۵۰} & \lr{۲۴٫۶} \\
\lr{P10} & \lr{۱۱۰} & \lr{۴۹} & \lr{۷۹} & \lr{۸۰} & \lr{۱۱۰} & \lr{۱۹٫۳} \\
\lr{P11} & \lr{۹۵} & \lr{۳۳} & \lr{۶۳} & \lr{۶۵} & \lr{۹۵} & \lr{۱۷٫۸} \\
\lr{P12} & \lr{۹۰} & \lr{۱۹} & \lr{۴۹} & \lr{۶۰} & \lr{۹۰} & \lr{۱۷٫۳} \\
\bottomrule
\end{tabular}
\end{table}


\fig{v2_history_P1.png}{0.85}{خطای \lr{validation} \lr{P1} در طول آموزش: چپ \lr{K}، راست \lr{E}+\lr{F}. \lr{K} در دوره‌ی ۸۰ بهترین بود و در دوره‌ی ۱۱۰ \lr{freeze} شد؛ \lr{EEG} و \lr{fusion} تا دوره‌ی ۱۵۰ ادامه دادند.}

\subsection{هر عضله و سقف داده}

جدول~\ref{tab:v2-musc} نشان می‌دهد \lr{K} در همه‌ی پنج عضله از مدل پایه بهتر است.

ستون «سقف» جدول می‌گوید در بهترین حالت چه عددی ممکن است. این عدد \lr{VAF} خود پوش واقعی پس از صاف کردن با \lr{low-pass} ۲ هرتز است؛ یعنی دقت یک پیش‌بین هموار که همه‌چیز را درست می‌داند و فقط نوسان‌های سریع را نمی‌تواند بسازد. \lr{FD} محدودترین عضله است: سقفش ۷۴ درصد است و فعالیتش تغییر کمی دارد (\lr{modulation depth} کم: صدک ۸۰ فقط حدود دو برابر صدک ۲۰ است). برای \lr{AD}، \lr{BR} و \lr{ED} فاصله‌ی \lr{K} تا سقف حدود ۸ تا ۱۱ واحد است؛ یعنی هنوز جا برای بهتر شدن هست.

\begin{table}[htbp]
\centering
\small
\caption{\lr{VAF} (درصد) برای هر عضله، میانگین روی ۱۲ شرکت‌کننده؛ \lr{smooth ceiling} = \lr{VAF} خودِ هدفِ هموارشده با \lr{low-pass} ۲ هرتز؛ \lr{modulation depth} = نسبت صدک ۸۰ به صدک ۲۰ هدف.}
\label{tab:v2-musc}
\begin{tabular}{ccccccc}
\toprule
عضله & پایه & \lr{K} & \lr{E} & \lr{F} & \lr{smooth ceiling} & \lr{modulation depth} \\
\midrule
\lr{AD} & \lr{۸۵٫۹} & \lr{۸۸٫۶} & \lr{۵۸٫۰} & \lr{۸۸٫۷} & \lr{۹۳٫۷} & \lr{۱۲٫۲} \\
\lr{BR} & \lr{۷۹٫۶} & \lr{۸۲٫۵} & \lr{۴۴٫۴} & \lr{۸۲٫۹} & \lr{۹۰٫۵} & \lr{۳٫۲} \\
\lr{FD} & \lr{۴۹٫۵} & \lr{۵۳٫۰} & \lr{۱۸٫۱} & \lr{۵۲٫۷} & \lr{۷۴٫۲} & \lr{۲٫۰} \\
\lr{ED} & \lr{۸۲٫۴} & \lr{۸۶٫۱} & \lr{۵۲٫۴} & \lr{۸۶٫۴} & \lr{۹۳٫۳} & \lr{۵٫۰} \\
\lr{FDI} & \lr{۶۳٫۷} & \lr{۶۸٫۶} & \lr{۲۲٫۷} & \lr{۶۸٫۵} & \lr{۸۷٫۱} & \lr{۳٫۵} \\
\bottomrule
\end{tabular}
\end{table}


\subsection{دامنه در هر وزن جسم}

یکی از کاستی‌های مدل پایه این بود که خطای دامنه به وزن جسم بستگی داشت. جدول~\ref{tab:v2-load} این را برای سه وزن مقایسه می‌کند. دو چیز سنجیده شده است: \lr{height ratio} در زمان نگه‌داشتن، و همبستگی سطح فعالیت بین بلندکردن‌هایی که وزن یکسان دارند.

جزء \lr{lift} در \lr{loss function} و ویژگی‌های تاریخچه دقیقاً برای همین خطا طراحی شدند و اثرشان روشن است. مدل پایه وزن سبک را \lr{۴٫۵} درصد و وزن متوسط را \lr{۵٫۵} درصد بیش‌برآورد می‌کرد. \lr{height ratio} \lr{K} در هر سه وزن بین \lr{۱٫۰۰۴} و \lr{۱٫۰۰۹} است، یعنی خطای کمتر از یک درصد. معیار پذیرش دوم این بود که این عدد بین \lr{۰٫۹۵} و \lr{۱٫۰۵} باشد.

همبستگی سطح فعالیت بین بلندکردن‌های هم‌وزن هم در هر سه وزن بالا رفته است؛ مثلاً در وزن سبک از \lr{۰٫۳۹} به \lr{۰٫۵۳}. یعنی \lr{K} بهتر از مدل پایه تشخیص می‌دهد کدام بلندکردن کمی بیشتر یا کمتر از بقیه زور می‌خواهد.

\begin{table}[htbp]
\centering
\small
\caption{سطح فعالیت در مرحله‌ی نگه‌داشتن، برای هر وزن جسم (میانگین روی شرکت‌کننده‌ها، سری‌ها و عضله‌ها). \lr{height ratio} = میانگین پیش‌بینی به میانگین هدف؛ مقدار مطلوب ۱.}
\label{tab:v2-load}
\begin{tabular}{ccccc}
\toprule
وزن جسم & \lr{height ratio} پایه & \lr{height ratio} \lr{K} & \lr{r} بین بلندکردن‌ها پایه & \lr{r} بین بلندکردن‌ها \lr{K} \\
\midrule
سبک & \lr{۱٫۰۴۵} & \lr{۱٫۰۰۵} & \lr{۰٫۳۹۱} & \lr{۰٫۵۲۹} \\
متوسط & \lr{۱٫۰۵۵} & \lr{۱٫۰۰۴} & \lr{۰٫۵۰۱} & \lr{۰٫۵۷۰} \\
سنگین & \lr{۱٫۰۲۴} & \lr{۱٫۰۰۹} & \lr{۰٫۳۹۴} & \lr{۰٫۴۶۱} \\
\bottomrule
\end{tabular}
\end{table}


\subsection{\lr{triptych figures}}

برای اینکه ببینیم مدل در عمل چه می‌کند، هر سری آزمون به تکه‌های ۱۰ ثانیه‌ای تقسیم و برای هر تکه یک شکل سه‌پنلی (\lr{triptych}) ساخته شد. در ادامه برای شرکت‌کننده‌های ۱ تا ۴ و برای هر سری آزمون یک شکل آمده است (ثانیه‌ی ۳۰ تا ۴۰ هر سری).

\begin{itemize}
\item \textbf{پنل ۱: پیش‌بینی در برابر واقعیت.} پوش واقعی (سیاه) کنار \lr{K} (نارنجی)، \lr{E} (آبی) و \lr{F} (سبز خط‌چین)، با واحد اصلی. خطوط عمودی رخدادهای آزمایش‌اند: روشن شدن چراغ، اولین تماس، جدا شدن جسم از سطح، گذاشتن و رها کردن. ناحیه‌ی زرد بازه‌ی نگه‌داشتن است. ستون راست $r$ همین تکه را می‌دهد.
\item \textbf{پنل ۲: اجزای خروجی.} $K_{base}$ (خاکستری)، $K=K_{base}\cdot g$ (نارنجی) و اصلاح \lr{EEG}، $F_{eeg}$ (بنفش)، در حوزه‌ی جذر. ستون راست بازه‌ی ضریب تاریخچه و نسبت اندازه‌ی اصلاح \lr{EEG} به \lr{K} را می‌دهد.
\item \textbf{پنل ۳: توجه بین عضله‌ها.} \lr{attention rollout} \lr{muscle graph} جریان سینماتیک، به‌صورت تغییر نسبت به میانگین استراحت همان سری. هر بلوک یک عضله‌ی گیرنده و هر سطر یک عضله‌ی فرستنده است.
\end{itemize}

الگوی همه‌ی شکل‌ها یکی است. \lr{F} دقیقاً روی \lr{K} افتاده و $F_{eeg}$ نزدیک صفر است. \lr{E} زمان‌بندی کلی هر بلندکردن را درست دنبال می‌کند، اما جهش‌های کوتاه و تیز پوش را ندارد.

\fig{trip_P1_S04.png}{0.82}{\lr{P1}، سری آزمون ۴، ثانیه‌ی ۳۰ تا ۴۰.}
\fig{trip_P1_S05.png}{0.82}{\lr{P1}، سری آزمون ۵، ثانیه‌ی ۳۰ تا ۴۰.}
\fig{trip_P2_S05.png}{0.82}{\lr{P2}، سری آزمون ۵، ثانیه‌ی ۳۰ تا ۴۰.}
\fig{trip_P2_S06.png}{0.82}{\lr{P2}، سری آزمون ۶، ثانیه‌ی ۳۰ تا ۴۰.}
\fig{trip_P3_S04.png}{0.82}{\lr{P3}، سری آزمون ۴، ثانیه‌ی ۳۰ تا ۴۰.}
\fig{trip_P3_S06.png}{0.82}{\lr{P3}، سری آزمون ۶، ثانیه‌ی ۳۰ تا ۴۰.}
\fig{trip_P4_S04.png}{0.82}{\lr{P4}، سری آزمون ۴، ثانیه‌ی ۳۰ تا ۴۰.}
\fig{trip_P4_S05.png}{0.82}{\lr{P4}، سری آزمون ۵، ثانیه‌ی ۳۰ تا ۴۰.}
\clearpage
